% Emacs, this is -*-latex-*-

\ex{Distribui\'c\~oes preditivas para modelos ocultos de Markov}
\label{ex:predictive-distributions-for-hidden-markov-models}
Para o modelo oculto de Markov
$$ p(h_{1:d},v_{1:d}) = p(v_1|h_1)p(h_1)\prod_{i=2}^d
p(v_i|h_i)p(h_i|h_{i-1})$$ assuma que voc\^e possui observa\'c\~oes para $v_i$,
$i=1, \ldots, u < d$.

\begin{exenumerate}
    \item \label{q:h-predictive-hmm-alt} Use passagem de mensagens para computar $p(h_t|v_{1:u})$ para $u<t\le
    d$. Para fins de concretude, voc\^e pode considerar o caso $d=6, u=2,
    t=4$.
    
    \begin{solution}
        O grafo de fatores para $d=6, u=2$, com mensagens que s\~ao necess\'arias
        para o c\'alculo de $p(h_t|v_{1:u})$ para $t=4$, \'e o seguinte:
        \begin{center}
            \scalebox{1}{
                \begin{tikzpicture}[ugraph,minimum size=1cm, inner sep=3pt]
                    
                    \node[fact, label=above:$\phi_1$] (f1) at (-1,2) {};;
                    \node[cont] (h1) at (0,2) {$h_1$};
                    \node[fact, label=above:$\phi_2$] (f2) at (1,2) {};
                    \node[cont] (h2) at (2,2) {$h_2$};
                    \node[fact, label=above: {\footnotesize$p(h_3 | h_2)$}] (f3) at (3,2) {};
                    \node[cont] (h3) at (4,2) {$h_3$};
                    \node[fact, label=above: {\footnotesize$p(h_4 | h_3)$}] (f4) at (5,2) {};
                    \node[cont] (h4) at (6,2) {$h_4$};
                    \node[fact, label=above:$p(h_5 | h_4)$] (f5) at (7,2) {};
                    \node[cont] (h5) at (8,2) {$h_5$};
                    \node[fact, label=above:$p(h_6 | h_5)$] (f6) at (9,2) {};
                    \node[cont] (h6) at (10,2) {$h_6$};
                    
                    \node[cont] (v3) at (4,0) {$v_3$};
                    \node[fact, label={[xshift=0.1cm,  yshift=0cm]left: {\footnotesize$p(v_3 | h_3)$}}] (fv3) at (4,1) {};
                    \node[cont] (v4) at (6,0) {$v_4$};
                    \node[fact, label={[xshift=0.1cm,  yshift=0cm]left: {\footnotesize$p(v_4 | h_4)$}}] (fv4) at (6,1) {};
                    \node[cont] (v5) at (8,0) {$v_5$};
                    \node[fact, label={[xshift=0.1cm,  yshift=0cm]left: {\footnotesize$p(v_5 | h_5)$}}] (fv5) at (8,1) {};
                    \node[cont] (v6) at (10,0) {$v_6$};
                    \node[fact, label={[xshift=0.1cm,  yshift=0cm]left: {\footnotesize$p(v_6 | h_6)$}}] (fv6) at (10,1) {};
                    
                    \draw(h5)--(v5);
                    \draw(h6)--(v6);
                    \draw(h1)--(h2);
                    \draw(h2)--(h3);
                    \draw(h3)--(h4);
                    \draw(h4)--(h5);
                    \draw(h5)--(h6);
                    \draw(f1) -- (h1);    
                    
                    \draw(f1)--(h1) node[midway, above, yshift=-7pt] {$\rightarrow$};
                    \draw(h1)--(f2) node[midway, above, yshift=-7pt] {$\rightarrow$};
                    \draw(f2)--(h2) node[midway, above, yshift=-7pt] {$\rightarrow$};
                    \draw(h2)--(f3) node[midway, above, yshift=-7pt] {$\rightarrow$};
                    \draw(f3)--(h3) node[midway, above, yshift=-7pt] {$\rightarrow$};
                    \draw(h3)--(f4) node[midway, above, yshift=-7pt] {$\rightarrow$};
                    \draw(f4)--(h4) node[midway, above, yshift=-7pt] {\red{$\rightarrow$}};
                    
                    \draw(h4)--(f5) node[midway, above, yshift=-7pt] {\red{$\leftarrow$}};
                    \draw(f5)--(h5) node[midway, above, yshift=-7pt] {$\leftarrow$};
                    \draw(h5)--(f6) node[midway, above, yshift=-7pt] {$\leftarrow$};
                    \draw(f6)--(h6) node[midway, above, yshift=-7pt] {$\leftarrow$};
                    
                    \draw(fv3)--(h3) node[midway, right, xshift=-7pt] {$\uparrow$};
                    \draw(v3)--(fv3) node[midway, right, xshift=-7pt] {$\uparrow$};
                    \draw(fv4)--(h4) node[midway, right, xshift=-7pt] {$\uparrow$};
                    \draw(v4)--(fv4) node[midway, right, xshift=-7pt] {$\uparrow$};
                    \draw(fv5)--(h5) node[midway, right, xshift=-7pt] {$\uparrow$};
                    \draw(v5)--(fv5) node[midway, right, xshift=-7pt] {$\uparrow$};
                    \draw(fv6)--(h6) node[midway, right, xshift=-7pt] {$\uparrow$};
                    \draw(v6)--(fv6) node[midway, right, xshift=-7pt] {$\uparrow$};
                    
                \end{tikzpicture}
            }
        \end{center}
        As mensagens das vari\'aveis vis\'iveis n\~ao observadas $v_i$ para seus
        correspondentes $h_i$, e.g.\ $v_3$ para $h_3$, s\~ao todas iguais a um. Al\'em disso,
        a mensagem do n\'o $p(h_5|h_4)$ para $h_4$ tamb\'em \'e igual a um. Isso
        ocorre porque todos os fatores envolvidos, $p(v_i|h_i)$ e
        $p(h_i|h_{i-1})$, somam um. Portanto, o grafo de fatores reduz-se a uma cadeia:
        \begin{center}
            \scalebox{1}{
                \begin{tikzpicture}[ugraph]          
                    \node[fact, label=above: $\phi_1$] (f1) at (-1.5,2) {};
                    \node[cont] (h1) at (0,2) {$h_1$};
                    \node[fact, label=above: $\phi_2$] (f2) at (1.5,2) {};
                    \node[cont] (h2) at (3,2) {$h_2$};
                    \node[fact, label=above: $p(h_3|h_2)$] (f3) at (4.5,2) {};
                    \node[cont] (h3) at (6,2) {$h_3$};
                    \node[fact, label=above: $p(h_4|h_3)$] (f4) at (7.5,2) {};
                    \node[cont] (h4) at (9,2) {$h_4$};
                    \draw(f1) -- (h1)  node[midway,above,sloped] {$\rightarrow$};
                    \draw(h1)--(f2);
                    \draw(f2)--(h2) node[midway,above,sloped] {$\rightarrow$};
                    \draw(h2)--(f3);
                    \draw(f3)--(h3) node[midway,above,sloped] {$\rightarrow$};
                    \draw(h3)--(f4);
                    \draw(f4)--(h4) node[midway,above,sloped] {\red{$\rightarrow$}}; 
            \end{tikzpicture}}
        \end{center}
        Como os n\'os de vari\'avel copiam as mensagens no caso de uma cadeia, mostramos apenas
        as mensagens de fator para vari\'avel.
        
        O grafo mostra que estamos essencialmente na mesma situa\'c\~ao que
        na filtragem, com a diferen\'ca de que usamos os fatores
        $p(h_s|h_{s-1})$ para $s\ge u+1$. Assim, podemos usar a filtragem para computar
        as mensagens at\'e o tempo $s=u$ e depois computar as demais
        mensagens com $p(h_s|h_{s-1})$ como fatores. Isso resulta no
        seguinte algoritmo:
        \begin{enumerate}
            \item Compute $\alpha(h_u)$ por filtragem.
            \item Para $s=u+1, \ldots, t$, compute
            \begin{equation}
                \alpha(h_s) = \sum_{h_{s-1}} p(h_s|h_{s-1}) \alpha(h_{s-1})
            \end{equation}
            \item A distribui\'c\~ao preditiva requerida \'e
            \begin{equation}
                p(h_t | v_{1:u}) = \frac{1}{Z} \alpha(h_t) \quad \quad Z = \sum_{h_t} \alpha(h_t)
            \end{equation}
        \end{enumerate}
        Para $s \ge u+1$, temos que
        \begin{align}
            \sum_{h_s} \alpha(h_s) &= \sum_{h_s} \sum_{h_{s-1}} p(h_s|h_{s-1}) \alpha(h_{s-1})\\
            & =  \sum_{h_{s-1}} \alpha(h_{s-1})
        \end{align}
        j\'a que $p(h_s|h_{s-1})$ \'e normalizada. Isso significa que a constante de normaliza\'c\~ao $Z$ acima \'e igual a
        \begin{equation}
            Z = \sum_{h_u} \alpha(h_u) = p(v_{1:u})
        \end{equation}
        que \'e a verossimilhan\'ca.
        
        Para a filtragem, vimos que $\alpha(h_s) \propto
        p(h_s|v_{1:s})$, $s\le u$. Os $\alpha(h_s)$ para todo $s>u$ s\~ao
        proporcionais a $p(h_s| v_{1:u})$. Isso pode ser visto notando que
        os argumentos acima valem para qualquer $t>u$. 
    \end{solution}
    
    \item Use passagem de mensagens para computar $p(v_t|v_{1:u})$ para $u<t\le
    d$. Para fins de concretude, voc\^e pode considerar o caso $d=6,
    u=2, t=4$.
    
    \begin{solution}
        O grafo de fatores para $d=6, u=2$, com mensagens que s\~ao necess\'arias
        para o c\'alculo de $p(v_t|v_{1:u})$ para $t=4$, \'e como segue:
        \begin{center}
            \scalebox{1}{
                \begin{tikzpicture}[ugraph,minimum size=1cm, inner sep=3pt]
                    
                    \node[fact, label=above:$\phi_1$] (f1) at (-1,2) {};;
                    \node[cont] (h1) at (0,2) {$h_1$};
                    \node[fact, label=above:$\phi_2$] (f2) at (1,2) {};
                    \node[cont] (h2) at (2,2) {$h_2$};
                    \node[fact, label=above: {\footnotesize$p(h_3 | h_2)$}] (f3) at (3,2) {};
                    \node[cont] (h3) at (4,2) {$h_3$};
                    \node[fact, label=above: {\footnotesize$p(h_4 | h_3)$}] (f4) at (5,2) {};
                    \node[cont] (h4) at (6,2) {$h_4$};
                    \node[fact, label=above:$p(h_5 | h_4)$] (f5) at (7,2) {};
                    \node[cont] (h5) at (8,2) {$h_5$};
                    \node[fact, label=above:$p(h_6 | h_5)$] (f6) at (9,2) {};
                    \node[cont] (h6) at (10,2) {$h_6$};
                    
                    \node[cont] (v3) at (4,0) {$v_3$};
                    \node[fact, label={[xshift=0.1cm,  yshift=0cm]left: {\footnotesize$p(v_3 | h_3)$}}] (fv3) at (4,1) {};
                    \node[cont] (v4) at (6,0) {$v_4$};
                    \node[fact, label={[xshift=0.1cm,  yshift=0cm]left: {\footnotesize$p(v_4 | h_4)$}}] (fv4) at (6,1) {};
                    \node[cont] (v5) at (8,0) {$v_5$};
                    \node[fact, label={[xshift=0.1cm,  yshift=0cm]left: {\footnotesize$p(v_5 | h_5)$}}] (fv5) at (8,1) {};
                    \node[cont] (v6) at (10,0) {$v_6$};
                    \node[fact, label={[xshift=0.1cm,  yshift=0cm]left: {\footnotesize$p(v_6 | h_6)$}}] (fv6) at (10,1) {};
                    
                    \draw(h5)--(v5);
                    \draw(h6)--(v6);
                    \draw(h1)--(h2);
                    \draw(h2)--(h3);
                    \draw(h3)--(h4);
                    \draw(h4)--(h5);
                    \draw(h5)--(h6);
                    \draw(f1) -- (h1);    
                    
                    \draw(f1)--(h1) node[midway, above, yshift=-7pt] {$\rightarrow$};
                    \draw(h1)--(f2) node[midway, above, yshift=-7pt] {$\rightarrow$};
                    \draw(f2)--(h2) node[midway, above, yshift=-7pt] {$\rightarrow$};
                    \draw(h2)--(f3) node[midway, above, yshift=-7pt] {$\rightarrow$};
                    \draw(f3)--(h3) node[midway, above, yshift=-7pt] {$\rightarrow$};
                    \draw(h3)--(f4) node[midway, above, yshift=-7pt] {$\rightarrow$};
                    \draw(f4)--(h4) node[midway, above, yshift=-7pt] {$\rightarrow$};
                    
                    \draw(h4)--(f5) node[midway, above, yshift=-7pt] {$\leftarrow$};
                    \draw(f5)--(h5) node[midway, above, yshift=-7pt] {$\leftarrow$};
                    \draw(h5)--(f6) node[midway, above, yshift=-7pt] {$\leftarrow$};
                    \draw(f6)--(h6) node[midway, above, yshift=-7pt] {$\leftarrow$};
                    
                    \draw(fv3)--(h3) node[midway, right, xshift=-7pt] {$\uparrow$};
                    \draw(v3)--(fv3) node[midway, right, xshift=-7pt] {$\uparrow$};
                    \draw(fv4)--(h4) node[midway, right, xshift=-7pt] {$\downarrow$};
                    \draw(v4)--(fv4) node[midway, right, xshift=-7pt] {\red{$\downarrow$}};
                    \draw(fv5)--(h5) node[midway, right, xshift=-7pt] {$\uparrow$};
                    \draw(v5)--(fv5) node[midway, right, xshift=-7pt] {$\uparrow$};
                    \draw(fv6)--(h6) node[midway, right, xshift=-7pt] {$\uparrow$};
                    \draw(v6)--(fv6) node[midway, right, xshift=-7pt] {$\uparrow$};
                    
                \end{tikzpicture}
            }
        \end{center}
        Devido aos fatores normalizados, como acima, as mensagens \`a direita
        de $h_t$ s\~ao todas iguais a um. Al\'em disso, as mensagens que sobem de
        $v_i$ para $h_i, i\neq t$, tamb\'em s\~ao todas um. Portanto, o grafo simplifica-se para uma cadeia.
        \begin{center}
            \scalebox{1}{
                \begin{tikzpicture}[ugraph]          
                    \node[fact, label=above: $\phi_1$] (f1) at (-1.5,2) {};
                    \node[cont] (h1) at (0,2) {$h_1$};
                    \node[fact, label=above: $\phi_2$] (f2) at (1.5,2) {};
                    \node[cont] (h2) at (3,2) {$h_2$};
                    \node[fact, label=above: $p(h_3|h_2)$] (f3) at (4.5,2) {};
                    \node[cont] (h3) at (6,2) {$h_3$};
                    \node[fact, label=above: $p(h_4|h_3)$] (f4) at (7.5,2) {};
                    \node[cont] (h4) at (9,2) {$h_4$};
                    \node[fact, label=above: $p(v_4|h_4)$] (fv4) at (10.5,2) {};
                    \node[cont] (v4) at (12,2) {$v_4$};
                    
                    \draw(f1) -- (h1)  node[midway,above,sloped] {$\rightarrow$};
                    \draw(h1)--(f2);
                    \draw(f2)--(h2) node[midway,above,sloped] {$\rightarrow$};
                    \draw(h2)--(f3);
                    \draw(f3)--(h3) node[midway,above,sloped] {$\rightarrow$};
                    \draw(h3)--(f4);
                    \draw(f4)--(h4) node[midway,above,sloped] {\blue{$\rightarrow$}};
                    \draw(h4)--(fv4);
                    \draw(fv4)--(v4) node[midway,above,sloped] {\red{$\rightarrow$}};
            \end{tikzpicture}}
        \end{center}
        A mensagem em azul \'e proporcional a $p(h_t|v_{1:u})$ computada
        na quest\~ao \ref{q:h-predictive-hmm-alt}. Assim, assuma que tenhamos
        computado $p(h_t|v_{1:u})$. A distribui\'c\~ao preditiva no
        n\'ivel das vis\'iveis \'e, portanto,
        \begin{equation}
            p(v_t | v_{1:u}) = \sum_{h_t} p(v_t|h_t) p(h_t|v_{1:u}).
        \end{equation}
        Isso decorre da passagem de mensagens, pois o \'ultimo n\'o ($h_4$ no
        grafo) apenas copia a mensagem (normalizada) e o pr\'oximo
        fator \'e igual a $p(v_t|h_t)$.
        
        Uma deriva\'c\~ao alternativa segue de defini\'c\~oes e
        opera\'c\~oes b\'asicas, juntamente com as independ\^encias em HMMs:
        \begin{align}
            &\text{\scriptsize (regra da soma)}&  p(v_t | v_{1:u}) &= \sum_{h_t} p(v_t, h_t|v_{1:u})\\
            &\text{\scriptsize (regra do produto)}&  & = \sum_{h_t} p(v_t|h_t, v_{1:u}) p(h_t|v_{1:u})\\
            &\text{\scriptsize ($v_t \independent v_{1:u} \mid h_t$)}&   & = \sum_{h_t} p(v_t|h_t) p(h_t|v_{1:u})
        \end{align}
    \end{solution}
    
\end{exenumerate}
% ---------------------------------------

\ex{Algoritmo de Viterbi}

Para o modelo oculto de Markov
$$ p(h_{1:t},v_{1:t}) = p(v_1|h_1)p(h_1)\prod_{i=2}^t
p(v_i|h_i)p(h_i|h_{i-1})$$ assuma que voc\^e possui observa\'c\~oes para $v_i$,
$i=1, \ldots, t$. Use o algoritmo max-sum para derivar um algoritmo iterativo para computar
\begin{equation}
    \hat{\h} = \argmax_{h_1, \ldots, h_t}  p(h_{1:t}|v_{1:t})
\end{equation}
Assuma que as vari\'aveis latentes $h_i$ possam assumir $K$ valores diferentes,
e.g.\ $h_i \in \{0,\ldots, K-1\}$. O algoritmo resultante \'e conhecido como
algoritmo de Viterbi.

\begin{solution}
    
    Primeiro formamos os fatores:
    \begin{align}
        \phi_1(h_1) &= p(v_1|h_1) p(h_1)& \phi_2(h_1, h_2) &= p(v_2|h_2) p(h_2|h_1)\\
        \ldots      &                   & \phi_t(h_{t-1},h_t) &= p(v_t|h_t) p(h_{t}|h_{t-1})
    \end{align}
    onde os $v_i$ s\~ao conhecidos e fixos. A posteriori $p(h_1, \ldots,
    h_t | v_1, \ldots, v_t)$ \'e ent\~ao representada pelo seguinte grafo de
    fatores (assumindo $t=4$).
    
    \begin{center}
        \scalebox{1}{
            \begin{tikzpicture}[ugraph]          
                \node[fact, label=above: $\phi_1$] (f1) at (-1.5,2) {};
                \node[cont] (h1) at (0,2) {$h_1$};
                \node[fact, label=above: $\phi_2$] (f2) at (1.5,2) {};
                \node[cont] (h2) at (3,2) {$h_2$};
                \node[fact, label=above: $\phi_3$] (f3) at (4.5,2) {};
                \node[cont] (h3) at (6,2) {$h_3$};
                \node[fact, label=above: $\phi_4$] (f4) at (7.5,2) {};
                \node[cont] (h4) at (9,2) {$h_4$};
                
                \draw(f1)--(h1);
                \draw(h1)--(f2);
                \draw(f2)--(h2);
                \draw(h2)--(f3);
                \draw(f3)--(h3);
                \draw(h3)--(f4);
                \draw(f4)--(h4);
        \end{tikzpicture}}
    \end{center}
    
    Para o algoritmo max-sum, escolhemos aqui $h_t$ como a raiz. 
    Iniciamos ent\~ao o algoritmo com $\mfxmess{\phi_1}{h_1}{1}(h_1) =
    \log \phi_1(h_1) = \log p(v_1|h_1) + \log p(h_1)$ e ent\~ao computamos as
    mensagens da esquerda para a direita, movendo-se da folha $\phi_1$ para a
    raiz $h_t$.
    
    Como estamos lidando com uma cadeia, os n\'os de vari\'avel, assim como no
    algoritmo sum-product, apenas copiam as mensagens de entrada. Portanto,
    basta computar as mensagens de fator para vari\'avel mostradas no
    grafo, e ent\~ao realizar o backtracking para $h_1$.
    
    \begin{center}
        \scalebox{1}{
            \begin{tikzpicture}[ugraph]          
                \node[fact, label=above: $\phi_1$] (f1) at (-1.5,2) {};
                \node[cont] (h1) at (0,2) {$h_1$};
                \node[fact, label=above: $\phi_2$] (f2) at (1.5,2) {};
                \node[cont] (h2) at (3,2) {$h_2$};
                \node[fact, label=above: $\phi_3$] (f3) at (4.5,2) {};
                \node[cont] (h3) at (6,2) {$h_3$};
                \node[fact, label=above: $\phi_4$] (f4) at (7.5,2) {};
                \node[cont] (h4) at (9,2) {$h_4$};
                
                \draw(f1) -- (h1)  node[midway,above,sloped] {$\rightarrow$};
                \draw(h1)--(f2);
                \draw(f2)--(h2) node[midway,above,sloped] {$\rightarrow$};
                \draw(h2)--(f3);
                \draw(f3)--(h3) node[midway,above,sloped] {$\rightarrow$};
                \draw(h3)--(f4);
                \draw(f4)--(h4) node[midway,above,sloped] {$\rightarrow$};
        \end{tikzpicture}}
    \end{center}
    
    Com $\mxfmess{h_{i-1}}{\phi_i}{1}(h_{i-1})=\mfxmess{\phi_{i-1}}{h_{i-1}}{1}(h_{i-1})$, a equa\'c\~ao de atualiza\'c\~ao fator-para-vari\'avel \'e:
    \begin{align}
        \mfxmess{\phi_i}{h_i}{1}(h_i) &= \max_{h_{i-1}} \{ \log \phi_i(h_{i-1}, h_i) + \mxfmess{h_{i-1}}{\phi_i}{1}(h_{i-1}) \}\\
        &= \max_{h_{i-1}} \{ \log \phi_i(h_{i-1}, h_i) + \mfxmess{\phi_{i-1}}{h_{i-1}}{1}(h_{i-1}) \}
    \end{align}
    Para simplificar a nota\'c\~ao, denote $\mfxmess{\phi_i}{h_i}{1}(h_i)$ por $V_i(h_i)$. Temos, assim:
    \begin{align}
        V_1(h_1) & = \log p(v_1|h_1) + \log p(h_1) \\
        V_i(h_i) & =  \max_{h_{i-1}} \{ \log \phi_i(h_{i-1}, h_i) + V_{i-1}(h_{i-1}) \} \quad \quad i=2, \ldots, t
    \end{align}
    Em geral, $V_1(h_1)$ e $V_i(h_i)$ s\~ao fun\'c\~oes que dependem de
    $h_1$ e $h_i$, respectivamente. Assumindo que os $h_i$ possam assumir
    os valores $0, \ldots, K-1$, as equa\'c\~oes acima podem ser escritas como:
    \begin{align}
        v_{1,k} &= \log p(v_1|k) + \log p(k) &  k&=0, \ldots, K-1\\
        v_{i,k} & =  \max_{m \in \{0, \ldots, K-1\}} \{ \log \phi_i(m, k) + v_{i-1,m} \} &  k&=0, \ldots, K-1, \quad i=2, \ldots, t,
    \end{align}
    Ao final do algoritmo, temos portanto uma matriz $\V$ de dimens\~ao $t \times K$
    com elementos $v_{i,k}$.
    
    A maximiza\'c\~ao pode ser realizada computando-se a matriz tempor\'aria
    $\A$ (via broadcasting) onde o elemento $(m,k)$ \'e $\log
    \phi_i(m, k) + v_{i-1,m}$. A maximiza\'c\~ao corresponde ent\~ao a
    determinar o valor m\'aximo em cada coluna.
    
    Para suportar o backtracking, quando computamos $V_i(h_i)$ ao maximizar
    sobre $h_{i-1}$, computamos simultaneamente a tabela de consulta:
    \begin{equation}
        \gamma_i^*(h_i)   =  \argmax_{h_{i-1}} \{ \log \phi_i(h_{i-1}, h_i) + V_{i-1}(h_{i-1}) \}
    \end{equation}
    Quando $h_i$ assume os valores $0, \ldots, K-1$, isso pode ser escrito como:
    \begin{equation}
        \gamma^*_{i,k} = \argmax_{m \in \{0, \ldots, K-1\}} \{ \log \phi_i(m, k) + v_{i-1,m} \}
    \end{equation}
    Este \'e o \'indice (da linha) do elemento m\'aximo em cada coluna da matriz tempor\'aria $\A$.
    
    Ap\'os computar $v_{t,k}$ e $\gamma^*_{t,k}$, realizamos ent\~ao o
    backtracking via:
    \begin{align}
        \hat{h}_t & = \argmax_{k}   v_{t,k}\\
        \hat{h}_i & = \gamma^*_{i+1,\hat{h}_{i+1}} \quad \quad i=t-1, \ldots, 1
    \end{align}
    Isso fornece recursivamente $\hat{\h} = (\hat{h}_1, \ldots, \hat{h}_t) = \argmax_{h_1, \ldots, h_t}  p(h_{1:t}|v_{1:t})$.
    
\end{solution}

% ---------------------------------------

\ex{Amostragem forward filtering backward sampling para modelos ocultos de Markov}
\label{ex:FFBS-alt}

Considere o modelo oculto de Markov especificado pelo seguinte DAG.

\begin{center}
    \scalebox{0.8}{
        \begin{tikzpicture}[dgraph, minimum size=0.5cm, inner sep=3pt]
            \node[cont] (h1) at (0,2) {$h_1$};
            \node       (aa) at (1,2) {};
            \node       (bb) at (2,2) {$\ldots$};
            \node       (cc) at (2,0) {$\ldots$};
            \node       (dd) at (3,2) {};
            
            \node[cont, inner sep=0em] (h2) at (5,2) {$h_{t-1}$};
            \node[cont] (h3) at (7,2) {$h_{t}$};
            
            \node       (a) at (8,2) {};
            \node       (b) at (9,2) {$\ldots$};
            \node       (c) at (9,0) {$\ldots$};
            \node       (d) at (10,2) {};
            
            \node[cont] (h4) at (12,2) {$h_n$};
            
            \node[cont] (v1) at (0,0) {$v_1$};
            \node[cont, inner sep=0em] (v2) at (5,0) {$v_{t-1}$};
            \node[cont] (v3) at (7,0) {$v_t$};
            \node[cont] (v4) at (12,0) {$v_n$};
            \draw(h1)--(v1);\draw(h2)--(v2);\draw(h3)--(v3);\draw(h4)--(v4);
            \draw(h1)--(aa);\draw(dd)--(h2);\draw(h2)--(h3);\draw(h3)--(a);\draw(d)--(h4);
    \end{tikzpicture}}
\end{center}
Assumimos que j\'a tenhamos executado a recurs\~ao alfa (filtragem) e
possamos computar $p(h_t | v_{1:t})$ para todo $t$. O objetivo \'e agora
gerar amostras de $p(h_1, \ldots, h_n | v_{1:n})$, i.e.\ trajet\'orias
inteiras $(h_1, \ldots, h_n)$ da posteriori. Note que isso n\~ao \'e o mesmo que amostrar das $n$ distribui\'c\~oes de filtragem $p(h_t |
v_{1:t})$. Al\'em disso, comparado ao algoritmo de Viterbi, a abordagem de amostragem
gera amostras da posteriori completa, em vez de apenas retornar
o estado mais prov\'avel e sua probabilidade correspondente.

\begin{exenumerate}
    \item Mostre que $p(h_1, \ldots, h_n|v_{1:n})$ forma uma
    cadeia de Markov de primeira ordem.
    
    \begin{solution}
        Existem v\'arias maneiras de mostrar isso. A mais simples \'e notar
        que o grafo n\~ao direcionado para o modelo oculto de Markov \'e o mesmo
        que o DAG, mas com as setas removidas, pois n\~ao h\'a colisores
        no DAG. Al\'em disso, o condicionamento corresponde \`a remo\'c\~ao de n\'os
        de um grafo n\~ao direcionado. Isso nos deixa com uma cadeia que
        conecta os $h_i$.
        \begin{center}
            \scalebox{0.8}{
                \begin{tikzpicture}[ugraph]
                    \node[cont] (h1) at (0,2) {$h_1$};
                    \node[cont] (h2) at (2,2) {$h_2$};
                    \node[cont] (h3) at (4,2) {$h_3$};
                    \node       (a) at (6,2) {};
                    \node       (b) at (7,2) {$\ldots$};
                    \node       (d) at (8,2) {};
                    \node[cont] (h4) at (10,2) {$h_n$};
                    \draw(h1)--(h2);\draw(h2)--(h3);\draw(h3)--(a);\draw(d)--(h4);
            \end{tikzpicture}}
        \end{center}
        
        Pela separa\'c\~ao de grafos, vemos que $p(h_1,\ldots, h_n|v_{1:n})$ forma
        uma cadeia de Markov de primeira ordem, de modo que e.g.\ $h_{1:t-1} \independent
        h_{t+1:n} | h_t$ (passado independente do futuro dado o
        presente).
        
    \end{solution}
    
    \item Como $p(h_1, \ldots, h_n|v_{1:n})$ \'e uma cadeia de Markov de
    primeira ordem, basta determinar $p(h_{t-1}|h_t, v_{1:n})$, a
    fun\'c\~ao de massa de probabilidade para $h_{t-1}$ dado $h_t$ e
    todos os dados $v_{1:n}$. Use passagem de mensagens para mostrar
    que
    \begin{equation}
        p(h_{t-1}, h_t| v_{1:n}) \propto \alpha(h_{t-1}) \beta(h_t) p(h_t|h_{t-1})p(v_t|h_t) 
    \end{equation}
    
    \begin{solution}
        Como todas as vis\'iveis est\~ao no conjunto de condicionamento, i.e.\ assumidas
        como observadas, podemos representar o modelo condicional $p(h_1, \ldots,
        h_n|v_{1:n})$ como uma \'arvore de fatores em cadeia, por exemplo, no caso
        de $n=4$:
        \begin{center}
            \scalebox{1}{
                \begin{tikzpicture}[ugraph]          
                    \node[fact, label=above: $\phi_1$] (f1) at (-1.5,2) {};
                    \node[cont] (h1) at (0,2) {$h_1$};
                    \node[fact, label=above: $\phi_2$] (f2) at (1.5,2) {};
                    \node[cont] (h2) at (3,2) {$h_2$};
                    \node[fact, label=above: $\phi_3$] (f3) at (4.5,2) {};
                    \node[cont] (h3) at (6,2) {$h_3$};
                    \node[fact, label=above: $\phi_4$] (f4) at (7.5,2) {};
                    \node[cont] (h4) at (9,2) {$h_4$};
                    \draw(f1) -- (h1);
                    \draw(h1)--(f2);
                    \draw(f2)--(h2);
                    \draw(h2)--(f3);
                    \draw(f3)--(h3);
                    \draw(h3)--(f4);
                    \draw(f4)--(h4);
            \end{tikzpicture}}
        \end{center}
        Combinando as distribui\'c\~oes de emiss\~ao $p(v_s|h_s)$ (e a marginal
        $p(h_1)$) com as distribui\'c\~oes de transi\'c\~ao $p(h_s|h_{s-1})$,
        obtemos os fatores:
        \begin{align}
            \phi_1(h_1) & = p(h_1) p(v_1|h_1)\\
            \phi_s(h_{s-1}, h_s) &= p(h_s|h_{s-1})p(v_s|h_s) \quad \text{para } t=2, \ldots, n
        \end{align}
        Vemos pela \'arvore de fatores que $h_{t-1}$ e $h_t$ s\~ao
        vizinhos, estando anexados ao mesmo n\'o de fator
        $\phi_t(h_{t-1}, h_t)$, por exemplo $\phi_3$ no caso de $p(h_{2}, h_3 | v_{1:4})$.
        
        Pelas regras de passagem de mensagens, a conjunta $p(h_{t-1}, h_{t} |
        v_{1:n})$ \'e proporcional a $\phi_t$ vezes as mensagens
        que entram em $\phi_t$. O grafo seguinte mostra as mensagens para o
        caso de $p(h_{2}, h_3 | v_{1:4})$.
        \begin{center}
            \scalebox{1}{
                \begin{tikzpicture}[ugraph]          
                    \node[fact, label=above: $\phi_1$] (f1) at (-1.5,2) {};
                    \node[cont] (h1) at (0,2) {$h_1$};
                    \node[fact, label=above: $\phi_2$] (f2) at (1.5,2) {};
                    \node[cont] (h2) at (3,2) {$h_2$};
                    \node[fact, label=above: $\phi_3$] (f3) at (4.5,2) {};
                    \node[cont] (h3) at (6,2) {$h_3$};
                    \node[fact, label=above: $\phi_4$] (f4) at (7.5,2) {};
                    \node[cont] (h4) at (9,2) {$h_4$};
                    \draw(f1) -- (h1);
                    \draw(h1)--(f2);
                    \draw(f2)--(h2);
                    \draw(h2)--(f3) node[midway,above,sloped] {$\rightarrow$};
                    \draw(f3)--(h3) node[midway,above,sloped] {$\leftarrow$}; 
                    \draw(h3)--(f4);
                    \draw(f4)--(h4);
            \end{tikzpicture}}
        \end{center}
        Como os n\'os de vari\'avel recebem apenas mensagens \'unicas de qualquer
        dire\'c\~ao, eles copiam as mensagens, de modo que as mensagens que entram em
        $\phi_t$ s\~ao dadas por $\alpha(h_{t-1})$ e $\beta(h_t)$, mostradas
        abaixo em vermelho e azul, respectivamente.
        \begin{center}
            \scalebox{1}{
                \begin{tikzpicture}[ugraph]          
                    \node[fact, label=above: $\phi_1$] (f1) at (-1.5,2) {};
                    \node[cont] (h1) at (0,2) {$h_1$};
                    \node[fact, label=above: $\phi_2$] (f2) at (1.5,2) {};
                    \node[cont] (h2) at (3,2) {$h_2$};
                    \node[fact, label=above: $\phi_3$] (f3) at (4.5,2) {};
                    \node[cont] (h3) at (6,2) {$h_3$};
                    \node[fact, label=above: $\phi_4$] (f4) at (7.5,2) {};
                    \node[cont] (h4) at (9,2) {$h_4$};
                    \draw(f1) -- (h1);
                    \draw(h1)--(f2);
                    \draw(f2)--(h2) node[midway,above,sloped] {\red{$\rightarrow$}};
                    \draw(h2)--(f3);
                    \draw(f3)--(h3);
                    \draw(h3)--(f4) node[midway,above,sloped] {\blue{$\leftarrow$}}; 
                    \draw(f4)--(h4);
            \end{tikzpicture}}
        \end{center}
        Portanto, 
        \begin{align}
            p(h_{t-1}, h_t| v_{1:n}) &\propto \alpha(h_{t-1}) \beta(h_t) \phi_t(h_{t-1}, h_t)\\
            &\propto \alpha(h_{t-1}) \beta(h_t) p(h_t|h_{t-1}) p(v_t|h_t)
        \end{align}
        que \'e o resultado que desejamos mostrar.
    \end{solution}
    
    
    \item Mostre que $p(h_{t-1}|h_t, v_{1:n}) = \frac{\alpha(h_{t-1})}{\alpha(h_t)} p(h_t|h_{t-1}) p(v_t|h_t)$.
    
    \begin{solution}
        A condicional $p(h_{t-1}|h_t, v_{1:n})$ pode ser escrita como a raz\~ao:
        \begin{equation}
            p(h_{t-1}|h_t, v_{1:n}) = \frac{p(h_{t-1}, h_t| v_{1:n})}{p(h_t|v_{1:n})}.
        \end{equation}
        Acima, mostramos que o numerador satisfaz:
        \begin{equation}
            p(h_{t-1}, h_t| v_{1:n}) \propto \alpha(h_{t-1}) \beta(h_t) p(h_t|h_{t-1}) p(v_t|h_t).
        \end{equation}
        O denominador $p(h_t|v_{1:n})$ \'e proporcional a $\alpha(h_t)
        \beta(h_t)$, j\'a que \'e a distribui\'c\~ao de suaviza\'c\~ao que pode ser
        determinada via recurs\~ao alfa-beta.
        
        Normalmente, precisar\'iamos somar as mensagens sobre todos os valores de
        $(h_{t-1},h_t)$ para encontrar a constante de normaliza\'c\~ao do
        numerador. Para o denominador, ter\'iamos que somar sobre todos os valores de
        $h_t$. A seguir, argumento qualitativamente que essa soma n\~ao
        \'e necess\'aria; as constantes de normaliza\'c\~ao s\~ao ambas iguais a
        $p(v_{1:n})$.
        
        Come\'camos com um grafo de fatores e fatores que representam a
        conjunta $p(h_{1:n}, v_{1:n})$. A condicional $p(h_{1:n} | v_{1:n})$ \'e igual a:
        \begin{equation}
            \label{eq:h-given-v-hmm-alt}
            p(h_{1:n} | v_{1:n}) = \frac{p(h_{1:n}, v_{1:n})}{p(v_{1:n})}
        \end{equation}
        A passagem de mensagens \'e elimina\'c\~ao de vari\'aveis. Portanto, ao computar
        $p(h_t|v_{1:n})$ como $\alpha(h_t) \beta(h_t)$ de um grafo de fatores para
        $p(h_{1:n}, v_{1:n})$, s\'o precisamos dividir por
        $p(v_{1:n})$ para a normaliza\'c\~ao; a soma expl\'icita sobre $h_t$ n\~ao \'e
        necess\'aria. Em outras palavras,
        \begin{equation}
            p(h_t|v_{1:n}) =  \frac{\alpha(h_t) \beta(h_t)}{p(v_{1:n})}.
        \end{equation}
        Similarmente, $p(h_{t-1}, h_t|v_{1:n})$ tamb\'em \'e obtida de
        \eqref{eq:h-given-v-hmm-alt} por marginaliza\'c\~ao/elimina\'c\~ao de vari\'aveis.
        Novamente, ao computar $p(h_{t-1}, h_t|v_{1:n})$ como
        $\alpha(h_{t-1}) \beta(h_t) p(h_t|h_{t-1}) p(v_t|h_t)$ de um
        grafo de fatores para $p(h_{1:n}, v_{1:n})$, n\~ao precisamos
        somar explicitamente sobre todos os valores de $h_t$ e $h_{t-1}$ para
        normaliza\'c\~ao. A defini\'c\~ao dos fatores no grafo de fatores
        juntamente com \eqref{eq:h-given-v-hmm-alt} mostra que podemos simplesmente
        dividir por $p(v_{1:n})$. Isso resulta em:
        \begin{align}
            p(h_{t-1}, h_t| v_{1:n}) & = \frac{1}{p(v_{1:n})} \alpha(h_{t-1}) \beta(h_t) p(h_t|h_{t-1}) p(v_t|h_t).
        \end{align}
        
        A condicional desejada \'e, portanto:
        \begin{align}
            p(h_{t-1} | h_t, v_{1:n}) & = \frac{p(h_{t-1}, h_t| v_{1:n})}{p(h_t|v_{1:n})} \\
            & = \frac{\alpha(h_{t-1}) \beta(h_t) p(h_t|h_{t-1}) p(v_t|h_t)}{\alpha(h_t) \beta(h_t)}\\
            & = \frac{\alpha(h_{t-1}) p(h_t|h_{t-1}) p(v_t|h_t)}{\alpha(h_t)}
        \end{align}
        que \'e o resultado que quer\'iamos mostrar. Note que
        $\beta(h_t)$ cancela-se e que $p(h_{t-1} | h_t, v_{1:n})$
        envolve apenas os $\alpha$'s, a distribui\'c\~ao de transi\'c\~ao (forward)
        $p(h_t|h_{t-1})$ e a distribui\'c\~ao de emiss\~ao no tempo $t$.
    \end{solution}
    
    Obtemos, assim, o seguinte algoritmo para gerar amostras de
    $p(h_1, \ldots, h_n|v_{1:n})$:
    \begin{enumerate}
        \item Execute a recurs\~ao alfa (filtragem) para determinar todos os
        $\alpha(h_t)$ para frente no tempo, de $t=1, \ldots, n$.
        \item Amostre $h_n$ de $p(h_n | v_{1:n}) \propto \alpha(h_n)$.
        \item Volte no tempo usando
        \begin{equation}
            p(h_{t-1}|h_t, v_{1:n}) = \frac{\alpha(h_{t-1})}{\alpha(h_t)} p(h_t|h_{t-1}) p(v_t|h_t)
        \end{equation}
        para gerar amostras de $h_{t-1}|h_t, v_{1:n}$ para $t=n, \ldots, 2$.
    \end{enumerate}
    Este algoritmo \'e conhecido como forward filtering backward sampling (FFBS). 
    
\end{exenumerate}


% ----------------------------------------

\ex{Exerc\'icio de predi\'c\~ao}

Considere um modelo oculto de Markov com tr\^es vis\'iveis $v_1, v_2, v_3$
e tr\^es vari\'aveis ocultas $h_1, h_2, h_3$, que podem ser representadas
com o seguinte grafo de fatores:

\begin{center}
    \scalebox{0.9}{
        \begin{tikzpicture}[ugraph,minimum size=1cm, inner sep=3pt]
            \node[cont] (v1) at (0,0) {$v_1$};
            \node[fact, label={[xshift=0.1cm,  yshift=0cm]left: $p(v_1 | h_1)$}] (fv1) at (0,1) {};
            \node[cont] (v2) at (3,0) {$v_2$};
            \node[fact, label={[xshift=0.1cm,  yshift=0cm]left: $p(v_2 | h_2)$}] (fv2) at (3,1) {};
            \node[cont] (v3) at (6,0) {$v_3$};
            \node[fact, label={[xshift=0.1cm,  yshift=0cm]left: $p(v_3 | h_3)$}] (fv3) at (6,1) {};
            
            \node[fact, label=left: $p(h_1)$] (f1) at (-1.5,2) {};
            \node[cont] (h1) at (0,2) {$h_1$};
            \node[fact, label={[label distance=-7pt]above: $p(h_2 | h_1)$}] (f2) at (1.5,2) {};
            \node[cont] (h2) at (3,2) {$h_2$};
            \node[fact, label={[label distance=-7pt]above: $p(h_3 | h_2)$}] (f3) at (4.5,2) {};
            \node[cont] (h3) at (6,2) {$h_3$};
            
            \draw(h1)--(v1);
            \draw(h2)--(v2);
            \draw(h3)--(v3);
            \draw(h1)--(h2);
            \draw(h2)--(h3);
            \draw(f1) -- (h1);
            
    \end{tikzpicture}}
\end{center}
Esta quest\~ao trata do c\'alculo da probabilidade preditiva $p(v_3=1 | v_1 =1)$.

\begin{exenumerate}
    
    \item O grafo de fatores abaixo representa $p(h_1, h_2, h_3, v_2, v_3
    \mid v_1 =1)$. Forne\'ca uma equa\'c\~ao que defina $\phi_A$ em termos dos
    fatores no grafo de fatores acima. 
    
    \begin{center}
        \scalebox{0.9}{
            \begin{tikzpicture}[ugraph,minimum size=1cm, inner sep=3pt]
                \node[cont] (v2) at (3,0) {$v_2$};
                \node[fact, label={[xshift=0.1cm,  yshift=0cm]left: $p(v_2 | h_2)$}] (fv2) at (3,1) {};
                \node[cont] (v3) at (6,0) {$v_3$};
                \node[fact, label={[xshift=0.1cm,  yshift=0cm]left: $p(v_3 | h_3)$}] (fv3) at (6,1) {};
                
                \node[cont] (h1) at (0,2) {$h_1$};
                \node[fact, label={[label distance=-7pt]above: $\phi_A$}] (f2) at (1.5,2) {};
                \node[cont] (h2) at (3,2) {$h_2$};
                \node[fact, label={[label distance=-7pt]above: $p(h_3 | h_2)$}] (f3) at (4.5,2) {};
                \node[cont] (h3) at (6,2) {$h_3$};
                
                \draw(h2)--(v2);
                \draw(h3)--(v3);
                \draw(h1)--(h2);
                \draw(h2)--(h3);
                
        \end{tikzpicture}}
    \end{center}
    
    \begin{solution}
        $\phi_A(h_1, h_2) \propto p(v_1 | h_1) p(h_1) p(h_2 | h_1)$ com $v_1=1$.
    \end{solution}
    
    \item Assuma, al\'em disso, que todas as vari\'aveis sejam
    bin\'arias, $h_i \in \{0,1\}$, $v_i \in \{0,1\}$; que $p(h_1=1)=0.5$, e que as distribui\'c\~oes de transi\'c\~ao e emiss\~ao sejam, para todo $i$, dadas por:
    
    \begin{center}
        \begin{tabular}{lll}
            \toprule
            $p(h_{i+1}|h_i)$ & $h_{i+1}$ &$h_i$ \\
            \midrule
            0 & 0 & 0\\
            1 & 1 & 0\\
            1 & 0 & 1\\
            0 & 1 & 1\\
            \bottomrule
        \end{tabular}
        \hspace{5ex}
        \begin{tabular}{lll}
            \toprule
            $p(v_i|h_i)$ & $v_i$ &$h_i$ \\
            \midrule
            0.6 & 0 & 0\\
            0.4 & 1 & 0\\
            0.4 & 0 & 1\\
            0.6 & 1 & 1\\
            \bottomrule
        \end{tabular}
    \end{center}
    Calcule os valores num\'ericos do fator $\phi_A$.
    
    \begin{solution}
        Dada a defini\'c\~ao das probabilidades de transi\'c\~ao e emiss\~ao, temos $\phi_A(h_1,h_2) = 0$ se $h_1=h_2$. Para $h_1=0, h_2 =1$, obtemos:
        \begin{align}
            \phi_A(h_1=0, h_2=1) &= p(v_1=1 | h_1=0) p(h_1=0) p(h_2=1 | h_1=0) \\
            & = 0.4 \cdot 0.5 \cdot 1\\
            & = \frac{4}{10} \cdot \frac{1}{2}\\
            & = \frac{2}{10} = 0.2
        \end{align}
        Para $h_1=1, h_2 =0$, obtemos:
        \begin{align}
            \phi_A(h_1=1, h_2=0) &= p(v_1=1 | h_1=1) p(h_1=1) p(h_2=0 | h_1=1) \\
            & = 0.6 \cdot 0.5 \cdot 1\\
            & = \frac{6}{10} \cdot \frac{1}{2}\\
            & = \frac{3}{10} = 0.3
        \end{align}
        Assim:
        \begin{center}
            \begin{tabular}{lll}
                \toprule
                $\phi_A(h_1, h_2)$ & $h_1$ &$h_2$ \\
                \midrule
                0 & 0 & 0\\
                0.3 & 1 & 0\\
                0.2 & 0 & 1\\
                0 & 1 & 1\\
                \bottomrule
            \end{tabular}
        \end{center}
    \end{solution}
    
    \item Denomine a mensagem do n\'o de vari\'avel $h_2$ para o n\'o de fator $p(h_3 |
    h_2)$ por $\alpha(h_2)$. Use passagem de mensagens para computar $\alpha(h_2)$
    para $h_2=0$ e $h_2=1$. Relate os valores de quaisquer mensagens intermedi\'arias
    que precisem ser computadas para o c\'alculo de $\alpha(h_2)$.
    
    \begin{solution}
        A mensagem de $h_1$ para $\phi_A$ \'e um. A mensagem de
        $\phi_A$ para $h_2$ \'e:
        \begin{align}
            \fxmess{\phi_A}{h_2}{1}(h_2=0) & = \sum_{h_1} \phi_A(h_1, h_2=0)\\
            & = 0.3\\
            \fxmess{\phi_A}{h_2}{1}(h_2=1) & = \sum_{h_1} \phi_A(h_1, h_2=1)\\
            & = 0.2
        \end{align}
        Como $v_2$ n\~ao \'e observado e $p(v_2|h_2)$ \'e normalizado, a mensagem de $p(v_2|h_2)$ para $h_2$ \'e igual a um.
        
        Isso significa que a mensagem de $h_2$ para $p(h_3 | h_2)$, que \'e
        $\alpha(h_2)$, \'e igual a $\fxmess{\phi_A}{h_2}{1}(h_2)$, i.e.\ 
        \begin{align}
            \alpha(h_2=0) & = 0.3\\
            \alpha(h_2=1) & = 0.2
        \end{align}
    \end{solution}
    
    \item Com $\alpha(h_2)$ definida como acima, use passagem de mensagens para mostrar
    que a probabilidade preditiva $p(v_3=1 | v_1=1)$ pode ser expressa em termos de $\alpha(h_2)$
    como:
    \begin{equation}
        p(v_3=1 | v_1=1) = \frac{ x \alpha(h_2=1) + y \alpha(h_2=0)}{\alpha(h_2=1)+\alpha(h_2=0)}
    \end{equation}
    e relate os valores de $x$ e $y$.  
    
    \begin{solution}
        Dada a defini\'c\~ao de $p(h_3| h_2)$, a mensagem $\fxmess{p(h_3|h_2)}{h_3}{1}(h_3)$ \'e:
        \begin{align}
            \fxmess{p(h_3|h_2)}{h_3}{1}(h_3=0) & = \alpha(h_2 = 1)\\
            \fxmess{p(h_3|h_2)}{h_3}{1}(h_3=1) & = \alpha(h_2 = 0)
        \end{align}
        O n\'o de vari\'avel $h_3$ copia a mensagem, de modo que temos:
        \begin{align}
            \fxmess{p(v_3|h_3)}{v_3}{1}(v_3=0) & = \sum_{h_3} p(v_3=0 | h_3)  \fxmess{p(h_3|h_2)}{h_3}{1}(h_3) \\
            & =  p(v_3=0 | h_3=0)  \alpha(h_2 = 1) + p(v_3=0 | h_3=1)  \alpha(h_2 = 0)\\
            & = 0.6  \alpha(h_2 = 1) + 0.4  \alpha(h_2 = 0)\\
            \fxmess{p(v_3|h_3)}{v_3}{1}(v_3=1) & = \sum_{h_3} p(v_3=1 | h_3)  \fxmess{p(h_3|h_2)}{h_3}{1}(h_3) \\
            & = p(v_3=1 | h_3=0)  \alpha(h_2 = 1) + p(v_3=1 | h_3=1)  \alpha(h_2 = 0)\\
            & = 0.4 \alpha(h_2=1) + 0.6 \alpha(h_2=0)
        \end{align}
        
        Temos, portanto:
        \begin{align}
            p(v_3 = 1| v_1=1) &=  \frac{0.4 \alpha(h_2=1) + 0.6 \alpha(h_2=0)}{ 0.4 \alpha(h_2=1) + 0.6 \alpha(h_2=0) +  0.6  \alpha(h_2 = 1) + 0.4  \alpha(h_2 = 0)}\\
            & =  \frac{0.4 \alpha(h_2=1) + 0.6 \alpha(h_2=0)}{\alpha(h_2=1)+\alpha(h_2=0)}
        \end{align}
        Os valores solicitados de $x$ e $y$ s\~ao: $x = 0.4$, $y = 0.6$. 
    \end{solution}
    
    \item Calcule o valor num\'erico de $p(v_3 =1 | v_1=1)$.
    
    \begin{solution}
        Inserindo os n\'umeros obtemos $\alpha(h_2=0)+\alpha(h_2=1) = 0.5 = 1/2$, de modo que:
        \begin{align}
            p(v_3 = 1| v_1=1) &= \frac{0.4 \cdot 0.2 + 0.6 \cdot 0.3}{\frac{1}{2}}\\
            &= 2 \cdot (0.08 + 0.18)\\
            &= 2 \cdot 0.26\\
            &= 0.52
        \end{align}
    \end{solution}
    
\end{exenumerate}


% ----------------------------------


\ex{Modelos ocultos de Markov e mudan\'ca de medida}

Adotamos aqui uma perspectiva de mudan\'ca de medida na recurs\~ao alfa.

Considere o seguinte grafo direcionado para um modelo oculto de Markov
onde os $y_i$ correspondem a vari\'aveis observadas (vis\'iveis) e os
$x_i$ a vari\'aveis n\~ao observadas (ocultas/latentes).

\begin{center}
    \scalebox{1}{
        \begin{tikzpicture}[dgraph]
            \node[cont] (x1) at (0,2) {$x_1$};
            \node[cont] (x2) at (2,2) {$x_2$};
            \node[cont] (x3) at (4,2) {$x_3$};
            \node       (a) at (6,2) {};
            \node       (b) at (7,2) {$\ldots$};
            \node       (c) at (7,0) {$\ldots$};
            \node       (d) at (8,2) {};
            \node[cont] (x4) at (10,2) {$x_n$};
            \node[cont] (y1) at (0,0) {$y_1$};
            \node[cont] (y2) at (2,0) {$y_2$};
            \node[cont] (y3) at (4,0) {$y_3$};
            \node[cont] (y4) at (10,0) {$y_n$};
            \draw(x1)--(y1);\draw(x2)--(y2);\draw(x3)--(y3);\draw(x4)--(y4);
            \draw(x1)--(x2);\draw(x2)--(x3);\draw(x3)--(a);\draw(d)--(x4);
    \end{tikzpicture}}
\end{center}

O modelo conjunto para $\x=(x_1, \ldots, x_n)$ e $\y = (y_1, \ldots,
y_n)$ \'e, portanto:
\begin{align}
    p(\x,\y) = p(x_1) \prod_{i=2}^n p(x_i|x_{i-1}) \prod_{i=1}^n p(y_i|x_i).
\end{align}


\begin{exenumerate}
    
    \item Mostre que
    \begin{equation}
        p(x_1, \ldots, x_n, y_1, \ldots, y_t) =  p(x_1) \prod_{i=2}^n p(x_i|x_{i-1}) \prod_{i=1}^t p(y_i|x_i)
    \end{equation}
    para $t=0, \ldots, n$. Consideramos o caso $t=0$ como correspondendo a $p(x_1, \ldots, x_n)$,
    \begin{equation}
        p(x_1, \ldots, x_n) =  p(x_1) \prod_{i=2}^n p(x_i|x_{i-1}).
    \end{equation}
    
    \begin{solution}
        O resultado segue integrando/somando $y_{t+1} \ldots n$.
        \begin{align}
            p(x_1, \ldots, x_n, y_1, \ldots, y_t) & = \int p(x_1, \ldots, x_n, y_1, \ldots, y_n) \ud y_{t+1} \ldots \ud y_n \\
            & = \int  p(x_1) \prod_{i=2}^n p(x_i|x_{i-1}) \prod_{i=1}^n p(y_i|x_i)\ud y_{t+1} \ldots \ud y_n \\
            & =  p(x_1) \prod_{i=2}^n p(x_i|x_{i-1}) \prod_{i=1}^t p(y_i|x_i) \int  \prod_{i=t+1}^n p(y_i|x_i)\ud y_{t+1} \ldots \ud y_n \\
            &=  p(x_1) \prod_{i=2}^n p(x_i|x_{i-1}) \prod_{i=1}^t p(y_i|x_i)   \prod_{i=t+1}^n \underbrace{\int p(y_i|x_i)\ud y_{i}}_{=1} \\
            & =  p(x_1) \prod_{i=2}^n p(x_i|x_{i-1}) \prod_{i=1}^t p(y_i|x_i)
        \end{align}
        O resultado para $p(x_1, \ldots, x_n)$ \'e obtido quando integramos todas as
        vari\'aveis $y$.
    \end{solution}
    
    \item Mostre que $p(x_1, \ldots, x_n | y_1, \ldots, y_t)$, $t=0,
    \ldots, n$, se fatoriza como
    \begin{equation}
        p(x_1, \ldots, x_n | y_1, \ldots, y_t) \propto  p(x_1) \prod_{i=2}^n p(x_i|x_{i-1}) \prod_{i=1}^t g_i(x_i)
        \label{eq:HMM-conditional-factorisation-alt}
    \end{equation}
    onde $g_i(x_i) = p(y_i|x_i)$ para um valor fixo de $y_i$, e que sua constante de normaliza\'c\~ao $Z_t$ \'e igual \`a verossimilhan\'ca $p(y_1, \ldots, y_t)$.
    
    \begin{solution}
        O resultado segue da defini\'c\~ao b\'asica da condicional:
        \begin{align}
            p(x_1, \ldots, x_n | y_1, \ldots, y_t) & = \frac{p(x_1, \ldots, x_n, y_1, \ldots, y_t)}{p(y_1, \ldots, y_t)}
        \end{align}
        juntamente com a express\~ao para $p(x_1, \ldots, x_n, y_1, \ldots, y_t)$ quando os $y_i$ s\~ao mantidos fixos.
    \end{solution}
    
    \item Denomine $ p(x_1, \ldots, x_n|y_1, \ldots, y_t)$ por $p_t(x_1,
    \ldots, x_n)$. O \'indice $t \le n$ indica o tempo da
    \'ultima vari\'avel $y$ em que estamos condicionando. Mostre a seguinte
    recurs\~ao para $1 \le t \le n$:
    \begin{align}
        p_{t-1}(x_1, \ldots, x_t) &=
        \begin{cases}
            p(x_1) & \text{se } t=1\\
            p_{t-1}(x_1, \ldots, x_{t-1}) p(x_t|x_{t-1}) & \text{caso contr\'ario}
            \label{eq:HMM-extension-alt}
        \end{cases} && \text{\small (extens\~ao)}\\
        p_{t}(x_1, \ldots, x_{t}) &= \frac{1}{Z_t} p_{t-1}(x_1, \ldots, x_t) g_t(x_t) \label{eq:HMM-change-of-measure-alt} && \text{\small (mudan\'ca de medida)}\\
        Z_t & = \int p_{t-1}(x_t) g_t(x_t) \ud x_t 
    \end{align}
    Ao iterar de $t=1$ at\'e $t=n$, podemos assim computar recursivamente
    $p(x_1, \ldots, x_n|y_1, \ldots, y_n)$, incluindo sua constante de
    normaliza\'c\~ao $Z_n$, que \'e igual \`a verossimilhan\'ca $ Z_n = p(y_1, \ldots, y_n)$. 
    
    \begin{solution}
        Come\'camos com \eqref{eq:HMM-conditional-factorisation-alt} que mostra
        que, por defini\'c\~ao de $p_t(x_1, \ldots, x_n)$, temos:
        \begin{align}
            p_t(x_1, \ldots, x_n) & = p(x_1, \ldots, x_n | y_1, \ldots, y_t)\\
            &\propto p(x_1) \prod_{i=2}^n p(x_i|x_{i-1}) \prod_{i=1}^t g_i(x_i)
            \label{eq:HMM-conditional-factorisation-2-alt} 
        \end{align}
        Para $t=1$, temos:
        \begin{equation}
            p_1(x_1, \ldots, x_n) \propto  p(x_1) \prod_{i=2}^n p(x_i|x_{i-1}) g_1(x_1)
        \end{equation}
        Integrar em rela\'c\~ao a $x_2, \ldots, x_n$ resulta em:
        \begin{align}
            p_1(x_1) & = \int p_1(x_1, \ldots, x_n) \ud x_2 \ldots \ud x_n\\
            &\propto \int  p(x_1) \prod_{i=2}^n p(x_i|x_{i-1}) g_1(x_1)\ud x_2 \ldots \ud x_n\\
            &\propto p(x_1) g_1(x_1) \int  \prod_{i=2}^n p(x_i|x_{i-1})\ud x_2 \ldots \ud x_n\\
            &\propto p(x_1) g_1(x_1) \underbrace{\prod_{i=2}^n \int p(x_i|x_{i-1})\ud x_i}_{=1}\\
            &\propto p(x_1) g_1(x_1)
        \end{align}
        
        A constante de normaliza\'c\~ao \'e:
        \begin{align}
            Z_1 & = \int  p(x_1) g_1(x_1) \ud x_1
        \end{align}
        Isso estabelece o resultado para $t=1$.
        
        De \eqref{eq:HMM-conditional-factorisation-alt}, temos ainda:
        \begin{align}
            p_{t-1}(x_1, \ldots, x_n) & = p(x_1, \ldots, x_n | y_1, \ldots, y_{t-1})\\
            &\propto p(x_1) \prod_{i=2}^n p(x_i|x_{i-1}) \prod_{i=1}^{t-1} g_i(x_i)
        \end{align}
        Integrar em rela\'c\~ao a $x_{t+1}, \ldots, x_n$, portanto, fornece:
        \begin{align}
            p_{t-1}(x_1, \ldots, x_t) & = \int p_{t-1}(x_1, \ldots, x_n) \ud x_{t+1} \ldots \ud x_n\\
            &\propto \int p(x_1) \prod_{i=2}^n p(x_i|x_{i-1}) \prod_{i=1}^{t-1} g_i(x_i)\ud x_{t+1} \ldots \ud x_n\\
            & \propto  p(x_1) \prod_{i=2}^t p(x_i|x_{i-1}) \prod_{i=1}^{t-1} g_i(x_i) \int \prod_{i=t+1}^n p(x_i|x_{i-1}) \ud x_{t+1} \ldots \ud x_n\\
            & \propto  p(x_1) \prod_{i=2}^t p(x_i|x_{i-1}) \prod_{i=1}^{t-1} g_i(x_i) \prod_{i=t+1}^n \int p(x_i|x_{i-1}) \ud x_i\\
            & \propto  p(x_1) \prod_{i=2}^t p(x_i|x_{i-1}) \prod_{i=1}^{t-1} g_i(x_i)
        \end{align}
        Notando que o produto sobre os $g_i$ n\~ao envolve $x_t$ e
        que $p(x_t | x_{t-1})$ \'e uma fdp, temos al\'em disso:
        \begin{align}
            p_{t-1}(x_1, \ldots, x_{t-1}) &=\int  p_{t-1}(x_1, \ldots, x_t) \ud x_t\\
            &  \propto p(x_1) \prod_{i=2}^{t-1} p(x_i|x_{i-1}) \prod_{i=1}^{t-1} g_i(x_i)
        \end{align}
        Logo:
        \begin{equation}
            p_{t-1}(x_1, \ldots, x_t) =  p_{t-1}(x_1, \ldots, x_{t-1}) p(x_t|x_{t-1})
        \end{equation}
        Note que podemos usar o sinal de igual pois $p(x_t|x_{t-1})$ \'e uma
        fdp e, portanto, integra um. Isso \'e s\^o vezes chamado de
        ``extens\~ao'', pois as entradas para $p_{t-1}$ s\~ao estendidas de
        $(x_1, \ldots, x_{t-1})$ para $x_1, \ldots, x_t$.
        
        De \eqref{eq:HMM-conditional-factorisation-2-alt}, temos ainda:
        \begin{align}
            p_t(x_1, \ldots, x_n) & \propto  p_{t-1}(x_1, \ldots, x_n) g_t(x_t)
        \end{align}
        Integrar em rela\'c\~ao a $x_{t+1}, \ldots, x_n$ fornece, ent\~ao:
        \begin{align}
            p_t(x_1, \ldots, x_t) & \propto  p_{t-1}(x_1, \ldots, x_t) g_t(x_t)
        \end{align}
        Esta \'e uma mudan\'ca de medida de $p_{t-1}(x_1, \ldots, x_t)$ para
        $p_t(x_1, \ldots, x_t)$. Note que $p_{t-1}(x_1, \ldots, x_t)$ envolve
        $g_i$, e portanto observa\'c\~oes $y_i$, apenas at\'e o tempo $t-1$. A
        mudan\'ca de medida multiplica o fator adicional
        $g_t(x_t)=p(y_t|x_t)$ e, assim, incorpora a observa\'c\~ao no tempo $t$
        ao modelo.
        
        A recurs\~ao declarada completa-se computando a constante de normaliza\'c\~ao $Z_t$ para
        $p_t(x_1, \ldots, x_t)$, que \'e igual a:
        \begin{align}
            Z_t & = \int  p_{t-1}(x_1, \ldots, x_t) g_t(x_t) \ud x_1, \ldots \ud x_t\\
            & = \int g_t(x_t) \left[\int p_{t-1}(x_1, \ldots, x_t) \ud x_1, \ldots \ud x_{t-1} \right] \ud x_t\\
            & = \int g_t(x_t) p_{t-1}(x_t) \ud x_t
        \end{align}
        
        Esta recurs\~ao, e algumas generaliza\'c\~oes leves, formam a base
        para o que \'e conhecido como a ``recurs\~ao forward'' em filtragem
        de part\'iculas e Monte Carlo sequencial. Uma excelente introdu\'c\~ao a
        esses t\'opicos \'e o livro \citep{Chopin2020}.
        
    \end{solution}
    
    \item Use a recurs\~ao acima para derivar a seguinte forma da recurs\~ao alfa: 
    \begin{align}
        p_{t-1}(x_{t-1}, x_t) &= p_{t-1}(x_{t-1}) p(x_t|x_{t-1}) && \text{\small (extens\~ao)} \\
        p_{t-1}(x_t) & = \int  p_{t-1}(x_{t-1}, x_t) \ud x_{t-1} &&\text{\small (marginaliza\'c\~ao)}\\
        p_{t}(x_t) &= \frac{1}{Z_t} p_{t-1}(x_t) g_t(x_t)  &&\text{\small (mudan\'ca de medida)}\\
        Z_t & = \int p_{t-1}(x_t) g_t(x_t) \ud x_t
    \end{align}
    com $p_0(x_1) = p(x_1)$.
    
    O termo $p_{t}(x_t)$ corresponde a $\alpha(x_t)$ da
    recurs\~ao alfa ap\'os a normaliza\'c\~ao. Al\'em disso, $p_{t-1}(x_t)$ \'e a
    distribui\'c\~ao preditiva para $x_t$ dadas as observa\'c\~oes at\'e o tempo
    $t-1$. Multiplicar $p_{t-1}(x_t)$ por $g_t(x_t)$ resulta no novo
    $\alpha(x_t)$. O termo $g_t(x_t) = p(y_t|x_t)$ \'e s\^o vezes chamado
    de termo de ``corre\'c\~ao''. Vemos aqui que a corre\'c\~ao tem o
    efeito de uma mudan\'ca de medida, transformando a distribui\'c\~ao preditiva
    $p_{t-1}(x_t)$ na distribui\'c\~ao de filtragem $p_t(x_t)$.
    
    \begin{solution}
        Seja $t>1$. Com \eqref{eq:HMM-extension-alt}, temos:
        \begin{align}
            p_{t-1}(x_{t-1}, x_t) & = \int p_{t-1}(x_1, \ldots, x_t) \ud x_1 \ldots \ud x_{t-2}\\
            &=  \int p_{t-1}(x_1, \ldots, x_{t-1}) p(x_t|x_{t-1})  \ud x_1 \ldots \ud x_{t-2}\\
            &=  p(x_t|x_{t-1})  \int p_{t-1}(x_1, \ldots, x_{t-1}) \ud x_1 \ldots \ud x_{t-2}\\
            & =  p(x_t|x_{t-1}) p_{t-1}(x_{t-1})
        \end{align}
        o que prova a ``extens\~ao''.
        
        Com \eqref{eq:HMM-change-of-measure-alt}, temos:
        \begin{align}
            p_t(x_t) & = \int  p_{t}(x_1, \ldots, x_{t}) \ud x_1, \ldots \ud x_{t-1}\\
            &=  \frac{1}{Z_t} \int p_{t-1}(x_1, \ldots, x_t) g_t(x_t) \ud x_1, \ldots \ud x_{t-1}\\
            & = \frac{1}{Z_t} g_t(x_t) \int p_{t-1}(x_1, \ldots, x_t)  \ud x_1, \ldots \ud x_{t-1}\\
            & = \frac{1}{Z_t} g_t(x_t) p_{t-1}(x_t)
        \end{align}
        o que prova a ``mudan\'ca de medida''. Al\'em disso, a constante de
        normaliza\'c\~ao $Z_t$ \'e a mesma de antes. Assim, completar a
        itera\'c\~ao at\'e $t=n$ produz a verossimilhan\'ca $p(y_1, \ldots, y_n) =
        Z_n$ como um subproduto da recurs\~ao. A inicializa\'c\~ao da
        recurs\~ao com $p_0(x_1)=p(x_1)$ tamb\'em \'e a mesma de cima.
        
    \end{solution}
    
\end{exenumerate}


% --------------------------------------------------------------


\ex{Filtragem de Kalman}
\label{ex:kalman-alt}

Consideramos aqui a filtragem para modelos ocultos de Markov com distribui\'c\~oes
de transi\'c\~ao e emiss\~ao Gaussianas. Por simplicidade, assumimos
vari\'aveis ocultas e observ\'aveis unidimensionais. Denotamos a
fun\'c\~ao de densidade de probabilidade de uma vari\'avel aleat\'oria Gaussiana $x$ com
m\'edia $\mu$ e vari\^ancia $\sigma^2$ por $\Gauss(x | \mu, \sigma^2)$:
\begin{equation}
    \Gauss(x | \mu, \sigma^2) = \frac{1}{\sqrt{2 \pi \sigma^2}} \exp\left[ -\frac{(x-\mu)^2}{2\sigma^2} \right].
\end{equation}
As distribui\'c\~oes de transi\'c\~ao e emiss\~ao s\~ao assumidas como:
\begin{align}
    p(h_s | h_{s-1}) &= \Gauss(h_s | A_s h_{s-1}, B^2_s) \label{eq:transition-alt}\\
    p(v_s | h_s) & = \Gauss(v_s | C_s h_s, D^2_s). \label{eq:emission-alt}
\end{align}
A distribui\'c\~ao $p(h_1)$ \'e assumida Gaussiana com par\^ametros conhecidos.
$A_s, B_s, C_s, D_s$ tamb\'em s\~ao assumidos conhecidos.


\begin{exenumerate}
    
    \item Mostre que $h_s$ e $v_s$, conforme definidos nas seguintes equa\'c\~oes de atualiza\'c\~ao
    e observa\'c\~ao
    \begin{align}
        h_s & = A_s h_{s-1} + B_s \xi_s \label{eq:transition-model-alt}\\
        v_s & = C_s h_s  + D_s \eta_s \label{eq:emission-model-alt}
    \end{align}
    seguem as distribui\'c\~oes condicionais em \eqref{eq:transition-alt} e
    \eqref{eq:emission-alt}. As vari\'aveis aleat\'orias $\xi_s$ e $\eta_s$ s\~ao
    independentes das outras vari\'aveis no modelo e seguem uma
    distribui\'c\~ao Gaussiana normal padr\~ao, e.g. $\xi_s \sim \Gauss(\xi_s |
    0, 1)$.\\ Dica: Para duas constantes $c_1$ e $c_2$, $y = c_1 + c_2 x$
    \'e Gaussiana se $x$ for Gaussiana. Em outras palavras, uma
    transforma\'c\~ao afim de uma Gaussiana \'e Gaussiana.
    
    As equa\'c\~oes significam que $h_s$ \'e obtido escalando $h_{s-1}$ e
    adicionando ru\'ido com vari\^ancia $B_s^2$. O valor observado $v_s$ \'e
    obtido escalando a oculta $h_s$ e corrompendo-a com
    ru\'ido de observa\'c\~ao Gaussiano de vari\^ancia $D_s^2$.
    
    \begin{solution}
        
        Por hip\'otese, $\xi_s$ \'e Gaussiana. Como condicionamos em $h_{s-1}$,
        $A_s h_{s-1}$ em \eqref{eq:transition-model-alt} \'e uma constante, e como
        $B_s$ tamb\'em \'e uma constante, $h_s$ \'e Gaussiana. 
        
        O que temos que mostrar em seguida \'e que \eqref{eq:transition-model-alt} define
        a mesma m\'edia e vari\^ancia condicionais que a Gaussiana condicional em
        \eqref{eq:transition-alt}: A esperan\'ca condicional de $h_s$ dado
        $h_{s-1}$ \'e:
        \begin{align}
            \E(h_s | h_{s-1}) & = A_s h_{s-1} + \E( B_s \xi_s) &&\text{(pois condicionamos em $h_{s-1}$)}\\
            & = A_s h_{s-1} + B_s \E(\xi_s) && \text{(por linearidade da esperan\'ca)}\\
            & = A_s h_{s-1} && \text{(pois $\xi_s$ tem m\'edia zero)}
        \end{align}
        A vari\^ancia condicional de $h_s$ dado $h_{s-1}$ \'e:
        \begin{align}
            \var(h_s | h_{s-1}) & =  \var( B_s \xi_s) &&\text{(pois condicionamos em $h_{s-1}$)}\\
            & = B_s^2 \var(\xi_s) && \text{(por propriedades da vari\^ancia)}\\
            & = B_s^2 && \text{(pois $\xi_s$ tem vari\^ancia um)}
        \end{align}
        Vemos que a m\'edia e a vari\^ancia condicionais de $h_s$ dado $h_{s-1}$ coincidem com aquelas em
        \eqref{eq:transition-alt}. E como $h_s$ dado $h_{s-1}$ \'e Gaussiana conforme argumentado acima,
        o resultado segue.
        
        Exatamente o mesmo racioc\'inio tamb\'em se aplica ao caso de
        \eqref{eq:emission-model-alt}. Condicional a $h_s$, $v_s$ \'e Gaussiana
        porque \'e uma transforma\'c\~ao afim de uma Gaussiana. A m\'edia condicional de $v_s$ dado $h_s$ \'e:
        \begin{align}
            \E(v_s | h_s) & = C_s h_s + \E( D_s \eta_s) &&\text{(pois condicionamos em $h_s$)}\\
            & = C_s h_s + D_s \E(\eta_s) && \text{(por linearidade da esperan\'ca)}\\
            & = C_s h_s && \text{(pois $\eta_s$ tem m\'edia zero)}
        \end{align}
        A vari\^ancia condicional de $v_s$ dado $h_s$ \'e:
        \begin{align}
            \var(v_s | h_s) & =  \var( D_s \eta_s) &&\text{(pois condicionamos em $h_s$)}\\
            & = D_s^2 \var(\eta_s) && \text{(por propriedades da vari\^ancia)}\\
            & = D_s^2 && \text{(pois $\eta_s$ tem vari\^ancia um)}
        \end{align}
        Portanto, condicional a $h_s$, $v_s$ \'e Gaussiana com m\'edia e vari\^ancia como em \eqref{eq:emission-alt}.
        
    \end{solution}
    
    
    \item Mostre que
    \begin{align}
        \int \Gauss(x | \mu, \sigma^2) \Gauss(y | Ax, B^2) \ud x & \propto \Gauss(y | A \mu, A^2 \sigma^2 + B^2) \label{eq:Gaussian-int-alt}
    \end{align} 
    Dica: Embora este resultado possa ser obtido por integra\'c\~ao, uma abordagem que evita isso \'e a seguinte: Primeiro note que $\Gauss(x | \mu, \sigma^2) \Gauss(y | Ax, B^2)$ \'e proporcional \`a fdp conjunta de $x$ e $y$. Podemos assim considerar que a integral corresponde ao c\'alculo da marginal de $y$ a partir da conjunta. Usar a equival\^encia das Equa\'c\~oes \eqref{eq:transition-alt}-\eqref{eq:emission-alt} e \eqref{eq:transition-model-alt}-\eqref{eq:emission-model-alt}, e o fato de que a soma ponderada de duas vari\'aveis aleat\'orias Gaussianas \'e uma vari\'avel aleat\'oria Gaussiana, permite ent\~ao obter o resultado.
    \begin{solution}
        Seguimos o procedimento descrito acima. As duas densidades Gaussianas correspondem \`as equa\'c\~oes:
        \begin{align}
            x & = \mu + \sigma \xi\\
            y &= A x + B \eta
        \end{align}
        onde $\xi$ e $\eta$ s\~ao vari\'aveis aleat\'orias normais padr\~ao independentes. A m\'edia de $y$ \'e:
        \begin{align}
            \E(y) & = A \E(x) + B \E(\eta)\\
            & = A \mu
        \end{align}
        onde usamos a linearidade da esperan\'ca e $\E(\eta)=0$. A vari\^ancia de $y$ \'e:
        \begin{align}
            \var(y) & = \var(A x) + \var (B \eta) \quad \quad \text{(j\'a que $x$ e $\eta$ s\~ao independentes)}\\
            & = A^2 \var(x) + B^2 \var(\eta) \quad \quad \text{(por propriedades da vari\^ancia)}\\
            & = A^2 \sigma^2 + B^2
        \end{align}
        Como $y$ \'e a soma (ponderada) de duas Gaussianas, ela pr\'opria \'e Gaussiana
        e, portanto, sua distribui\'c\~ao \'e completamente definida por sua
        m\'edia e vari\^ancia, de modo que:
        \begin{align}
            y & \sim \Gauss(y | A \mu, A^2 \sigma^2+B^2).
        \end{align}
        Agora, o produto $\Gauss(x | \mu, \sigma^2) \Gauss(y | Ax, B^2)$ \'e proporcional \`a fdp conjunta de $x$ e $y$, de modo que a integral pode ser considerada como correspondendo \`a marginaliza\'c\~ao de $x$, e portanto seu resultado \'e proporcional \`a densidade de $y$, que \'e $\Gauss(y | A \mu, A^2 \sigma^2+B^2)$.
    \end{solution}
    
    \item Mostre que
    \begin{equation}
        \Gauss(x | m_1, \sigma_1^2) \Gauss(x | m_2, \sigma_2^2) \propto \Gauss(x | m_3, \sigma_3^2) \label{eq:Gaussian-mult-alt}
    \end{equation}
    onde
    \begin{align}
        \sigma^2_3 & = \left( \frac{1}{\sigma_1^2} + \frac{1}{\sigma_2^2}\right)^{-1} = \frac{\sigma_1^2 \sigma_2^2}{\sigma_1^2+\sigma_2^2}\\
        m_3 & = \sigma_3^2\left( \frac{m_1}{\sigma_1^2}+ \frac{m_2}{\sigma_2^2}\right) = m_1 + \frac{\sigma_1^2}{\sigma_1^2+\sigma_2^2}(m_2-m_1)
    \end{align}
    \emph{Dica: Trabalhe no dom\'inio do logaritmo negativo.}
    
    \begin{solution}
        Mostramos o resultado usando uma t\'ecnica cl\'assica chamada ``completar o quadrado'', veja e.g.\ \url{https://en.wikipedia.org/wiki/Completing_the_square}.
        
        Trabalhamos no dom\'inio do logaritmo (negativo) e usamos que:
        \begin{align}
            -\log \left[\Gauss(x | m, \sigma^2)\right] & = \frac{ (x-m)^2 }{2 \sigma^2} + \text{const} \\
            & = \frac{ x^2}{2\sigma^2} - x \frac{m}{\sigma^2} + \frac{m^2}{2\sigma^2} + \text{const} \\
            & = \frac{ x^2}{2\sigma^2} - x \frac{m}{\sigma^2} +  \text{const} \label{eq:expansion-alt}
        \end{align}
        onde const indica termos que n\~ao dependem de $x$. Obtemos assim:
        \begin{align}
            -\log \left[\Gauss(x | m_1, \sigma_1^2) \Gauss(x | m_2, \sigma_2^2) \right]  & = -\log \left[\Gauss(x | m_1, \sigma_1^2)\right]-\log \left[ \Gauss(x | m_2, \sigma_2^2) \right]\\
            &= \frac{ (x-m_1)^2 }{2 \sigma_1^2} + \frac{ (x-m_2)^2 }{2 \sigma_2^2} + \text{const}\\
            &= \frac{ x^2}{2\sigma_1^2} - x \frac{m_1}{\sigma_1^2} + \frac{ x^2}{2\sigma_2^2} - x \frac{m_2}{\sigma_2^2} + \text{const}\\
            &= \frac{x^2}{2} \left(\frac{1}{\sigma_1^2}+\frac{1}{\sigma_2^2}\right) -x \left( \frac{m_1}{\sigma_1^2}+\frac{m_2}{\sigma_2^2}\right) + \text{const}\\
            &=  \frac{x^2}{2 \sigma_3^2} - \frac{x}{\sigma_3^2} \sigma_3^2\left( \frac{m_1}{\sigma_1^2}+\frac{m_2}{\sigma_2^2}\right) + \text{const},
        \end{align}
        onde:
        \begin{align}
            \frac{1}{\sigma_3^2} &= \frac{1}{\sigma_1^2}+\frac{1}{\sigma_2^2}.
        \end{align}
        A compara\'c\~ao com \eqref{eq:expansion-alt} mostra que podemos escrever ainda:
        \begin{align}
            \frac{x^2}{2 \sigma_3^2} - \frac{x}{\sigma_3^2} \sigma_3^2\left( \frac{m_1}{\sigma_1^2}+\frac{m_2}{\sigma_2^2}\right) &=  \frac{ (x-m_3)^2 }{2 \sigma_3^2} + \text{const}
        \end{align}
        onde:
        \begin{align}
            m_3 & =  \sigma_3^2\left( \frac{m_1}{\sigma_1^2}+\frac{m_2}{\sigma_2^2}\right)
        \end{align}
        de modo que:
        \begin{align}
            -\log \left[\Gauss(x | m_1, \sigma_1^2) \Gauss(x | m_2, \sigma_2^2) \right]  & = \frac{ (x-m_3)^2 }{2 \sigma_3^2} + \text{const}
        \end{align}
        e assim:
        \begin{align}
            \Gauss(x | m_1, \sigma_1^2) \Gauss(x | m_2, \sigma_2^2) & \propto  \Gauss(x | m_3, \sigma_3^2).
        \end{align}
        Note que a identidade:
        \begin{align}
            m_3 & = \sigma_3^2\left( \frac{m_1}{\sigma_1^2}+ \frac{m_2}{\sigma_2^2}\right) = m_1 + \frac{\sigma_1^2}{\sigma_1^2+\sigma_2^2}(m_2-m_1)
        \end{align}
        \'e obtida como segue:
        \begin{align}
            \sigma_3^2\left( \frac{m_1}{\sigma_1^2}+ \frac{m_2}{\sigma_2^2}\right) & =  \frac{\sigma_1^2 \sigma_2^2}{\sigma_1^2+\sigma_2^2}\left( \frac{m_1}{\sigma_1^2}+ \frac{m_2}{\sigma_2^2}\right)\\
            & =  m_1 \frac{\sigma_2^2}{\sigma_1^2+\sigma_2^2}+ m_2 \frac{\sigma_1^2}{\sigma_1^2+\sigma_2^2}\\
            & = m_1 \left(1- \frac{\sigma_1^2}{\sigma_1^2+\sigma_2^2}\right) + m_2 \frac{\sigma_1^2}{\sigma_1^2+\sigma_2^2}\\
            & = m_1 + \frac{\sigma_1^2}{\sigma_1^2+\sigma_2^2} (m_2-m_1)
        \end{align}
        
    \end{solution}
    
    
    \item Podemos usar a ``recurs\~ao alfa'' para computar recursivamente $p(h_t | v_{1:t}) \propto
    \alpha(h_t)$ da seguinte forma:
    \begin{align}
        \alpha(h_1) & = p(h_1) \cdot  p(v_1 | h_1) &
        \alpha(h_s) &=  p(v_s | h_s) \sum_{h_{s-1}} p(h_s | h_{s-1}) \alpha(h_{s-1}).
    \end{align}
    Para vari\'aveis aleat\'orias cont\'inuas, a soma acima torna-se uma integral, de modo que:
    \begin{align}
        \alpha(h_s) &=  p(v_s | h_s) \int p(h_s | h_{s-1}) \alpha(h_{s-1}) \ud h_{s-1}.
        \label{eq:cont-alpha-alt}
    \end{align}
    
    Para refer\^encia, denotemos a integral por $I(h_s)$:
    \begin{equation}
        I(h_s) = \int p(h_s | h_{s-1}) \alpha(h_{s-1}) \ud h_{s-1}.
        \label{eq:I-def-kalman-alt}
    \end{equation}
    Note que $I(h_s)$ \'e proporcional \`a distribui\'c\~ao preditiva $p(h_s | v_{1:s-1})$.
    
    Para uma distribui\'c\~ao a priori Gaussiana para $h_1$ e probabilidade de emiss\~ao
    Gaussiana $p(v_1 | h_1)$, $\alpha(h_1) = p(h_1) \cdot p(v_1 | h_1)
    \propto p(h_1 | v_1)$ \'e proporcional a uma Gaussiana. Denotamos sua
    m\'edia por $\mu_1$ e sua vari\^ancia por $\sigma_1^2$, de modo que:
    \begin{equation}
        \alpha(h_1) \propto \Gauss(h_1 | \mu_1, \sigma_1^2).
    \end{equation}
    Assumindo que $\alpha(h_{s-1}) \propto \Gauss(h_{s-1} | \mu_{s-1}, \sigma_{s-1}^2)$ (o que vale para $s=2$), use a Equa\'c\~ao \eqref{eq:Gaussian-int-alt} para mostrar que:
    \begin{align}
        I(h_s) & \propto \Gauss(h_s | A_s \mu_{s-1}, P_s)
    \end{align}
    onde:
    \begin{align}
        P_s & = A^2_s \sigma^2_{s-1} + B_s^2.
    \end{align}
    
    \begin{solution}
        Podemos definir $\alpha(h_{s-1}) \propto \Gauss(h_{s-1} | \mu_{s-1}, \sigma_{s-1}^2)$. Como $p(h_s | h_{s-1})$ \'e Gaussiana, veja a Equa\'c\~ao \eqref{eq:transition-alt}, a Equa\'c\~ao \eqref{eq:I-def-kalman-alt} torna-se:
        \begin{align}
            I(h_s) &\propto \int  \Gauss(h_s | A_s h_{s-1}, B^2_s)  \Gauss(h_{s-1} | \mu_{s-1}, \sigma_{s-1}^2) \ud h_{s-1}.
        \end{align}
        A Equa\'c\~ao \eqref{eq:Gaussian-int-alt} com $x \equiv h_{s-1}$ e $y \equiv h_s$ fornece o resultado desejado:
        \begin{align}
            I(h_s) &\propto  \Gauss(h_s | A_s \mu_{s-1}, A_s^2 \sigma_{s-1}^2+B_s^2).
            \label{eq:prediction-Gaussian-alt}
        \end{align}
        Podemos entender a equa\'c\~ao da seguinte forma: Para calcular a
        m\'edia preditiva de $h_s $ dado $v_{1:s-1}$, propagamos a
        m\'edia de $h_{s-1} | v_{1:s-1}$ usando a equa\'c\~ao de atualiza\'c\~ao
        \eqref{eq:transition-model-alt}. Isso fornece o termo de m\'edia $A_s
        \mu_{s-1}$. Como $h_{s-1} | v_{1:s-1}$ tem vari\^ancia
        $\sigma_{s-1}^2$, a vari\^ancia de $h_s | v_{1:s-1}$ \'e dada por
        $A_s^2 \sigma_{s-1}^2$ mais um termo adicional, $B_s^2$, devido ao
        ru\'ido na propaga\'c\~ao para frente. Isso resulta no termo de vari\^ancia
        $A_s^2 \sigma_{s-1}^2+B_s^2$.
        
    \end{solution}
    
    \item Use a Equa\'c\~ao \eqref{eq:Gaussian-mult-alt} para mostrar que:
    \begin{align}
        \alpha(h_s) & \propto \Gauss\left(h_s | \mu_s , \sigma_s^2\right)
    \end{align}
    onde:
    \begin{align}
        \mu_s & =  A_s \mu_{s-1} + \frac{P_sC_s}{C_s^2 P_s+ D_s^2}\left(v_s - C_s A_s \mu_{s-1}\right)\\
        \sigma^2_s & =  \frac{P_s D_s^2}{P_s C_s^2+ D_s^2}
    \end{align}
    
    \begin{solution}
        Tendo computado $I(h_s)$, o passo final na recurs\~ao alfa \'e:
        \begin{align}
            \alpha(h_s) & = p(v_s | h_s) I(h_s)
        \end{align}
        Com a Equa\'c\~ao \eqref{eq:emission-alt} obtemos:
        \begin{align}
            \alpha(h_s) & \propto \Gauss(v_s | C_s h_s, D^2_s)  \Gauss(h_s | A_s \mu_{s-1}, P_s).
        \end{align}
        Notamos ainda que:
        \begin{align}
            \Gauss(v_s | C_s h_s, D_s^2) \propto \Gauss \left(h_s | C_s^{-1}
            v_s, \frac{D_s^2}{C_s^2}\right)
        \end{align}
        de modo que podemos aplicar a Equa\'c\~ao \eqref{eq:Gaussian-mult-alt} (com $m_1 = A \mu_{s-1}$, $\sigma_1^2 = P_s$):
        \begin{align}
            \alpha(h_s) & \propto  \Gauss\left(h_s | C_s^{-1} v_s, \frac{D_s^2}{C_s^2}\right) \Gauss(h_s | A_s \mu_{s-1}, P_s)\\
            & \propto \Gauss\left(h_s, \mu_s, \sigma_s^2\right)
        \end{align}
        com:
        \begin{align}
            \mu_s & =  A_s \mu_{s-1} + \frac{P_s}{P_s + \frac{D_s^2}{C_s^2}}\left(C_s^{-1} v_s - A_s \mu_{s-1}\right)\\
            & =  A_s \mu_{s-1} + \frac{P_sC_s^2}{C_s^2 P_s+ D_s^2}\left(C_s^{-1}v_s - A_s \mu_{s-1}\right)\\
            & =  A_s \mu_{s-1} + \frac{P_sC_s}{C_s^2 P_s+ D_s^2}\left(v_s - C_s A_s \mu_{s-1}\right)\\
            \sigma^2_s & = \frac{P_s \frac{D_s^2}{C_s^2}}{P_s + \frac{D_s^2}{C_s^2}}\\
            & =  \frac{P_s D_s^2}{P_s C_s^2+ D_s^2}\\
        \end{align}
    \end{solution}
    
    \item Mostre que $\alpha(h_s)$ pode ser reescrita como:
    \begin{align}
        \alpha(h_s) & \propto \Gauss\left(h_s | \mu_s , \sigma_s^2\right)
    \end{align}
    onde:
    \begin{align}
        \mu_s & = A_s \mu_{s-1} + K_s \left(v_s - C_s A_s \mu_{s-1}\right) \label{eq:mu-def-alt}\\
        \sigma_s^2 & = (1-K_sC_s)P_s\\
        K_s & = \frac{P_s C_s}{C_s^2 P_s + D_s^2}
    \end{align}
    Estas s\~ao as equa\'c\~oes do filtro de Kalman e $K_s$ \'e chamado de ganho do filtro de Kalman.
    
    \begin{solution}
        Come\'camos a partir de:
        \begin{align}
            \mu_s & = A_s \mu_{s-1} + \frac{P_sC_s}{C_s^2 P_s+ D_s^2}\left(v_s - C_s A_s \mu_{s-1}\right),
        \end{align}
        e vemos que:
        \begin{equation}
            \frac{P_sC_s}{C_s^2 P_s+ D_s^2} = K_s
        \end{equation}
        de modo que:
        \begin{align}
            \mu_s & = A_s \mu_{s-1} + K_s \left(v_s - C_s A_s \mu_{s-1}\right).
        \end{align}
        Para a vari\^ancia $\sigma_s^2$, temos:
        \begin{align}
            \sigma^2_s  & =  \frac{P_s D_s^2}{P_s C_s^2+ D_s^2}\\
            & =\frac{D_s^2}{P_s C_s^2+ D_s^2}  P_s \\
            & = \left(1- \frac{P_sC_s^2}{P_s C_s^2+ D_s^2}\right) P_s\\
            & = (1- K_s C_s) P_s,
        \end{align}
        que \'e o resultado desejado.
        
        O resultado da filtragem generaliza-se para latentes e
        vis\'iveis com valores vetoriais, onde as distribui\'c\~oes de transi\'c\~ao e
        emiss\~ao em \eqref{eq:transition-alt} e \eqref{eq:emission-alt} tornam-se:
        \begin{align}
            p(\h_s | \h_{s-1}) & = \Gauss(\h_s | \A_s \h_{s-1}, \Sigmab^h),\\
            p(\v_s | \h_s) & = \Gauss( \v_s | \C_s \h_s, \Sigmab^v),
        \end{align}
        onde $\Gauss()$ denota fdps Gaussianas multivariadas, e.g.:
        \begin{align}
            \Gauss(\v_s | \C_s \h_s, \Sigmab^v) &= \frac{1}{|\det( 2\pi \Sigmab^v)|^{1/2}}\exp\left( -\frac{1}{2} (\v_s-\C_s \h_s)^\top (\Sigmab^v)^{-1}(\v_s-\C_s \h_s)\right).
        \end{align}
        Temos ent\~ao:
        \begin{align}
            p(\h_t | \v_{1:t}) & = \Gauss( \h_t | \mub_t, \Sigmab_t)
        \end{align}
        onde a m\'edia e a vari\^ancia posteriores s\~ao computadas recursivamente como:
        \begin{align}
            \mub_s & = \A_s \mub_{s-1} + \K_s ( \v_s-\C_s \A_s \mub_{s-1} )\\
            \Sigmab_s & = (\I - \K_s \C_s) \PP_s\\
            \PP_s & = \A_s \Sigmab_{s-1} \A_s^\top + \Sigmab^h\\
            \K_s & = \PP_s \C_s^\top \left( \C_s \PP_s \C_s^\top + \Sigmab^v \right)^{-1}
        \end{align}
        e inicializadas com $\mub_1$ e $\Sigmab_1$ iguais \`a m\'edia
        e vari\^ancia de $p(\h_1 | \v_1)$. A matriz $\K_s$ \'e ent\~ao
        chamada de matriz de ganho de Kalman.
        
        O filtro de Kalman \'e amplamente aplic\'avel, veja
        e.g.\ \url{https://en.wikipedia.org/wiki/Kalman_filter}, e
        desempenhou um papel em eventos hist\'oricos, como o pouso
        na lua, veja e.g.\ \citep{Grewal2010}.
        
        Um exemplo da aplica\'c\~ao do filtro de Kalman ao rastreamento
        (tracking) \'e mostrado na Figura \ref{fig:Kalman-tracking-alt}.
        \begin{figure}[h]
            \centering
            \includegraphics[width=0.5 \textwidth]{Figure13-22}
            \caption{ \label{fig:Kalman-tracking-alt} Filtragem de Kalman para o rastreamento de um objeto em movimento. Os pontos azuis indicam as posi\'c\~oes reais do objeto em um espa\'co bidimensional em passos de tempo sucessivos, os pontos verdes denotam medi\'c\~oes ruidosas das posi\'c\~oes, e as cruzes vermelhas indicam as m\'edias das distribui\'c\~oes posteriores inferidas das posi\'c\~oes obtidas pela execu\'c\~ao das equa\'c\~oes de filtragem de Kalman. As covari\^ancias das posi\'c\~oes inferidas s\~ao indicadas pelas elipses vermelhas, que correspondem a contornos com um desvio padr\~ao. \citep[Figura 13.22]{Bishop2006}}
        \end{figure}
        
        
    \end{solution}
    
    \item Explique a Equa\'c\~ao \eqref{eq:mu-def-alt} em termos n\~ao t\'ecnicos. O que acontece se a vari\^ancia $D_s^2$ do ru\'ido de observa\'c\~ao tender a zero?
    
    \begin{solution}
        J\'a vimos que $A_s \mu_{s-1}$ \'e a m\'edia preditiva de $h_s$ dado $v_{1:s-1}$. O termo $C_s A_s \mu_{s-1}$
        \'e, portanto, a m\'edia preditiva de $v_s$ dadas as observa\'c\~oes at\'e
        ent\~ao, $v_{1:s-1}$. A diferen\'ca $v_s - C_s A_s \mu_{s-1}$ \'e,
        assim, o erro de predi\'c\~ao da observ\'avel. Como $\alpha(h_s)$ \'e
        proporcional a $p(h_s | v_{1:s})$ e $\mu_s$ sua m\'edia,
        vemos assim que a m\'edia posterior de $h_s | v_{1:s}$ \'e igual \`a
        m\'edia posterior de $h_s | v_{1:s-1}$, $A_s \mu_{s-1}$, atualizada
        pelo erro de predi\'c\~ao da observ\'avel ponderado pelo
        ganho de Kalman.
        
        Para $D_s^2 \to 0$, $K_s \to C_s^{-1}$ e:
        \begin{align}
            \mu_s & = A_s \mu_{s-1} + K_s \left(v_s - C_s A_s \mu_{s-1}\right) \\
            & = A_s \mu_{s-1} + C_s^{-1} \left(v_s - C_s A_s \mu_{s-1}\right) \\
            & = A_s \mu_{s-1} + C_s^{-1} v_s - A_s \mu_{s-1}\\
            & =  C_s^{-1} v_s,
        \end{align}
        de modo que a m\'edia posterior de $p(h_s | v_{1:s})$ \'e obtida
        invertendo a equa\'c\~ao de observa\'c\~ao. Al\'em disso, a vari\^ancia
        $\sigma_s^2$ de $h_s | v_{1:s}$ tende a zero, de modo que o valor
        de $h_s$ \'e conhecido com precis\~ao e \'e igual a $C_s^{-1} v_s$.
        
    \end{solution}
\end{exenumerate}